% Lösungen Einheit 4
\einheitenkopf{Einheit 4 von 4 · Sachaufgaben}
\zweigzeile{Hier lernst du, zu Zahlenrätseln und Rechtecken eine quadratische Gleichung aufzustellen, zu lösen und die Lösung auszuwählen, die zur Frage passt · neu in diesem Jahr (an der Oberschule je nach Buch Klasse 9 oder 10) · keine P10-Aufgabe · baut auf: Wurzelziehen (Einheit 1), Satz vom Nullprodukt (Einheit 2), p-q-Formel (Einheit 3)}
\erg{36}{a) $x_1 = 4$ \quad b) $x_1 = 8$ \quad c) beide \quad d) $t_2 = 2{,}5$ \quad e) beide ($5$ und $6$ oder $-6$ und $-5$)}
\erg{37}{a) $20$ oder $-20$ \quad b) $30$ oder $-30$ \quad c) $0{,}5$ oder $-0{,}5$ \quad d) $50$ oder $-50$ \quad e) $x^2 + 19 = 100$, $x^2 = 81$; die Zahl ist $9$ oder $-9$ \quad f) $x\cdot(x + 1) = 56$, $x^2 + x - 56 = 0$, $-0{,}5 \pm 7{,}5$; $x_1 = 7$, $x_2 = -8$; die Zahlen sind $7$ und $8$ oder $-8$ und $-7$}
\erg{38}{a) Breite $x$, Länge $x + 5$: $x\cdot(x + 5) = 84$, $x^2 + 5x - 84 = 0$, $-2{,}5 \pm \sqrt{6{,}25 + 84} = -2{,}5 \pm 9{,}5$; $x_1 = 7$, $x_2 = -12$ \quad b) $-12$ passt nicht: eine Breite ist nie negativ \quad c) Das Beet ist $7\,\text{m}$ breit und $12\,\text{m}$ lang. \quad d) $12\,\text{m}\cdot 7\,\text{m} = 84\,\text{m}^2$ \quad e) $(x + 2)\cdot(x + 5) = 88$, $x^2 + 7x - 78 = 0$, $-3{,}5 \pm \sqrt{12{,}25 + 78} = -3{,}5 \pm 9{,}5$; $x_1 = 6$, $x_2 = -13$; die Seite war $6\,\text{cm}$ lang}
\erg{39}{a) Fehler: die negative Lösung steht in der Antwort – eine Länge ist nie negativ; richtig: Der Teich ist $10\,\text{m}$ lang. \quad b) $(x + 6)^2 = 400$; $x + 6 = 20$ oder $x + 6 = -20$; $x_1 = 14$, $x_2 = -26$; eine Seite ist $14\,\text{m}$ lang}
\erg{40}{a) eine Zahl darf negativ sein: $14$ und $-14$ haben beide das Quadrat $196$; eine Seitenlänge ist positiv, also nur $14\,\text{cm}$ \quad b) nein: beim Zahlenrätsel (oder bei Temperaturen) passt auch eine negative Lösung; nur bei Längen, Flächen, Anzahlen und Zeiten nach dem Start fällt sie weg}
